\documentclass[a4paper]{article}

\usepackage{graphicx}
\usepackage{amsmath,amssymb}
\usepackage{minted}
\usepackage{multirow}
\usepackage{float}
\usepackage{todonotes}
\usepackage{forest}
\usepackage{xstring}
\usepackage{xcolor}
\usepackage[most]{tcolorbox}
\usepackage{xparse}
\usepackage{natbib}

\usetikzlibrary{graphs,graphdrawing}
\usegdlibrary{layered}

% for external graphviz
\input{/home/donkarlo/Dropbox/repo/nd_ai_project/src/nd_ai/knoweldge_representation/language/formal/computer/programming/tex/latex/code_clips/external_graphviz.tex}

% for tags
\input{/home/donkarlo/Dropbox/repo/nd_ai_project/src/nd_ai/knoweldge_representation/language/formal/computer/programming/tex/latex/parts/control_sequence/kinds/semtag/semtag.tex}

% for links
\input{/home/donkarlo/Dropbox/repo/nd_ai_project/src/nd_ai/knoweldge_representation/language/formal/computer/programming/tex/latex/code_clips/seehere.tex}

\title{Language related self organization}
\date{}

\begin{document}
    \maketitle
    \cite{witzany-2014-biological-self-organization} argues that “biological self-organization” should not be explained merely through mechanical cause-and-effect models. The author argues that in living beings, organization does not occur only through the physical exchange of matter, but rather depends on communicative processes; that is, cells, tissues, organs, and organisms interact with one another in a “meaningful” way through signals, and without these interactions biological coordination is not possible.
    
    The core of the paper is that biological communication does not consist only of a sender and a receiver, but, like natural language, has three layers: syntax, meaning, and context-dependent use. In other words, a biological signal does not produce only one fixed and mechanical reaction; the same signal can take on different meanings depending on the context, time, place, and condition of the living being. According to the author, it is precisely this dependence on context that separates life from non-living processes.
    
    Then the paper moves into genetics and says that DNA should not be regarded merely as a passive storehouse of information. The author emphasizes the role of non-coding RNAs, mobile genetic elements, and viruses, and says that these have active roles in the regulation, alteration, and rearrangement of the genome. In this view, biological novelty is not only the result of “random mutation,” but is to some extent the result of the action of agents that can manipulate and re-regulate the genetic text.
    
    Another important part of the paper is its critique of mathematical and formal views of language. The author says that models which see language or the genetic code only as a formal syntax or a computable structure cannot explain the main feature of natural language: namely, that meaning is formed within social context and real use. He extends this point to biology and says that biological codes, too, are not understood only through formal rules, but depend on living agents and the practical context in which signs are used.
    
    The overall conclusion of the paper is that biological self-organization is fundamentally a “communicative-semiotic” process, not merely a physical-mechanical one. That is, life should be understood as a network of living agents that coordinate themselves through signs, rules, and context-dependent interpretation. For this reason, the author defends a non-reductionist and non-mechanistic description of life.
    
    If I want to put it very simply, the paper’s point is this:
    A living being does not merely “react”; it “interprets signs,” and this interpretation is the basis of the organization of life.
    \bibliographystyle{plainnat}
    \bibliography{/home/donkarlo/Dropbox/repo/nd_shared_research/bibliography/bibliography.bib}
\end{document}